% Standalone problem document, generated from the adjacent README.md.
% Compile with XeLaTeX. Mathematical content is maintained in Markdown.
\documentclass[11pt,a4paper]{article}
\usepackage[fontsize=10.5pt]{scrextend}
\usepackage[margin=22mm,headheight=15pt,headsep=9mm,footskip=11mm]{geometry}
\usepackage{fontspec}
\setmainfont{texgyrepagella}[Extension=.otf,UprightFont=*-regular,BoldFont=*-bold,ItalicFont=*-italic,BoldItalicFont=*-bolditalic]
\setsansfont{texgyreheros}[Extension=.otf,UprightFont=*-regular,BoldFont=*-bold,ItalicFont=*-italic,BoldItalicFont=*-bolditalic]
\setmonofont{texgyrecursor}[Extension=.otf,UprightFont=*-regular,BoldFont=*-bold,ItalicFont=*-italic,BoldItalicFont=*-bolditalic,Scale=MatchLowercase]
\usepackage{amsmath,amssymb}
\usepackage{unicode-math}
\setmathfont{texgyrepagella-math.otf}
\usepackage{xcolor,microtype,enumitem,needspace,fancyhdr}
\usepackage{bookmark}
\definecolor{ink}{HTML}{17344B}
\definecolor{accent}{HTML}{197D80}
\definecolor{muted}{HTML}{596773}
\hypersetup{colorlinks=true,linkcolor=accent,urlcolor=accent,pdftitle={MF-12 - Realizing
arbitrary polynomial growth exponents by finite matrix
families},pdfauthor={Open Problems in Numerical Linear Algebra}}
\urlstyle{same}
\setlength{\parindent}{0pt}
\setlength{\parskip}{5pt plus 1pt minus 1pt}
\linespread{1.04}
\setlength{\emergencystretch}{3em}
\widowpenalty=10000
\clubpenalty=10000
\displaywidowpenalty=10000
\allowdisplaybreaks[1]
\setlist{nosep,leftmargin=1.5em}
\providecommand{\tightlist}{\setlength{\itemsep}{0pt}\setlength{\parskip}{0pt}}
\providecommand{\pandocbounded}[1]{#1}
\pagestyle{fancy}
\fancyhf{}
\fancyhead[L]{\sffamily\footnotesize\color{muted}OPEN PROBLEMS IN NUMERICAL LINEAR ALGEBRA}
\fancyhead[R]{\sffamily\bfseries\footnotesize\color{accent}MF-12}
\fancyfoot[L]{\parbox[b]{0.9\textwidth}{\sffamily\fontsize{8}{10}\selectfont\color{muted}Editorial ratings; a bounded search cannot certify that a problem remains unsolved.\\Literature check: 2026-09-22}}
\fancyfoot[R]{\sffamily\footnotesize\color{muted}\thepage}
\renewcommand{\headrulewidth}{0.3pt}
\renewcommand{\headrule}{\hbox to\headwidth{\color{accent}\leaders\hrule height \headrulewidth\hfill}}
\makeatletter
\renewcommand{\section}{\Needspace{5\baselineskip}\@startsection{section}{1}{0pt}{12pt plus 2pt minus 2pt}{4pt}{\sffamily\bfseries\large\color{ink}}}
\renewcommand{\subsection}{\Needspace{4\baselineskip}\@startsection{subsection}{2}{0pt}{9pt plus 1pt minus 1pt}{3pt}{\sffamily\bfseries\normalsize\color{ink}}}
\makeatother
\setcounter{secnumdepth}{0}
\begin{document}
{\sffamily\small\color{muted}Matrix functions and stability\par}
\vspace{3pt}
{\raggedright\hyphenpenalty=10000\exhyphenpenalty=10000\sffamily\bfseries\fontsize{21}{24}\selectfont\color{ink}Realizing
arbitrary polynomial growth exponents by finite matrix families\par}
\vspace{7pt}
{\sffamily\small\textbf{Difficulty:} challenging\quad\textbf{Importance:} interesting
to the community\par}
{\sffamily\small\bfseries\color{accent}Status: Lean verified\par}
\vspace{4pt}
{\color{accent}\hrule height 0.8pt}
\vspace{6pt}
\textbf{Rating rationale:} Challenging because finite families must
realize every exponent at every length; community impact connects
switched dynamics and asymptotic matrix-product growth.

\subsection{Lean verification --- 22 September
2026}\label{lean-verification-22-september-2026}

The complete original target is formally proved: every real exponent α≥0
is realized by two fixed distinct real matrices in a fixed positive
dimension. The actual maximal Euclidean operator norm has positive
two-sided polynomial bounds at every positive length, and its actual
nth-root limit is one. All four exports passed LeanCert kernel audits,
two independent nonauthor final reviews, and
\href{https://github.com/sidneyholden1/OpenProblemsInNLA/actions/runs/35696246513}{isolated
Linux Comparator/default-kernel verification} with rejection and sandbox
controls.

\href{https://github.com/sidneyholden1/OpenProblemsInNLA/blob/main/matrix-functions-and-stability/MF-12/lean/README.md}{Proof
and reviews} ·
\href{https://github.com/sidneyholden1/OpenProblemsInNLA/tree/1231b92f2b0235559e5fe5162f6b76d21249ffd0/matrix-functions-and-stability/MF-12/lean}{Immutable
proof revision} ·
\href{https://github.com/sidneyholden1/OpenProblemsInNLA/blob/main/docs/lean/verification-2026-09-22/MF-12/README.md}{Retained
evidence}. Mathematics: Matthew J. Colbrook, with Varney--Morris
antecedents retained. Formalization: Sidney Holden with OpenAI Codex
assistance. The proof covers the full canonical target; optional
optimal-dimension and rational-entry refinements are not formalized. No
novelty, source-author endorsement or external human review is claimed.

\subsection{Resolution --- 2026-09-11}\label{resolution-2026-09-11}

\textbf{Affirmative resolution.} Matthew J. Colbrook's
\href{https://github.com/ajt60gaibb/OpenProblemsInNLA/blob/main/references/colbrook-jsr-growth-2026-09-11/manuscripts/arbitrary_growth_exponents.pdf}{complete
manuscript, Theorem 1 and its proof in §§2--5} constructs two distinct
real matrices for every real \(\alpha\ge0\), with joint spectral radius
one and maximal length-\(k\) product norm between positive constant
multiples of \(k^\alpha\) for every integer \(k\ge1\). The matrices and
dimension are fixed once the exponent is chosen. Six-dimensional pairs
suffice for \(0<\alpha<1\); every nonnegative rational exponent can be
realized with dyadic-rational entries. The proof covers all switching
words for the upper bound and every length for the lower bound,
including small lengths. It asserts comparability, not convergence of
the normalized growth sequence.

The complete original proof passed
\href{https://github.com/ajt60gaibb/OpenProblemsInNLA/blob/main/references/colbrook-jsr-growth-2026-09-11/verification/reviews/MF-12-review.md}{independent
Codex-agent review}.
\href{https://github.com/ajt60gaibb/OpenProblemsInNLA/blob/main/references/colbrook-jsr-growth-2026-09-11/manuscripts/arbitrary_growth_exponents.tex}{Authored
TeX} ·
\href{https://github.com/ajt60gaibb/OpenProblemsInNLA/blob/main/references/colbrook-jsr-growth-2026-09-11/README.md}{Submission,
authorship and verification record}. The proof was developed with AI
assistance; no external human peer review or formal verification is
claimed. The original statement and prior evidence below are retained,
and the ratings above are historical. This entry no longer contributes
to the open count.

\subsection{Context and notation}\label{context-and-notation}

The joint spectral radius of a nonempty compact set
\(\mathcal M\subset\mathbb C^{d\times d}\) is

\[\widehat\rho(\mathcal M)=\lim_{k\to\infty}
\max_{A_1,\ldots,A_k\in\mathcal M}\|A_k\cdots A_1\|_2^{1/k}.\]

This definition also applies to finite real matrix sets. The ordinary
spectral radius of one matrix is written \(\rho(A)\).

For a compact nonempty \(\mathcal M\subset\mathbb R^{d\times d}\) with
\(\widehat\rho(\mathcal M)=1\), define its maximal product norm at
length \(k\) by

\[g_{\mathcal M}(k)=\max_{A_1,\ldots,A_k\in\mathcal M}\|A_k\cdots A_1\|_2.\]

This problem concerns growth within one family over time.
\href{https://github.com/ajt60gaibb/OpenProblemsInNLA/blob/main/matrix-functions-and-stability/MF-07/README.md}{MF-07}
instead requests a dimension-dependent bound uniform across families.

\subsection{Problem statement}\label{problem-statement}

For every real \(\alpha\ge0\), do there exist a positive integer \(d\),
a finite nonempty \(\mathcal M\subset\mathbb R^{d\times d}\) with
\(\widehat\rho(\mathcal M)=1\), and constants \(0< c\le C<\infty\) such
that

\[c k^\alpha\le g_{\mathcal M}(k)\le C k^\alpha
\qquad\text{for every integer }k\ge1?\]

\subsection{Reference and status
evidence}\label{reference-and-status-evidence}

Varney and Morris, \href{https://arxiv.org/html/2209.00449}{On marginal
growth rates of matrix products}, §7, Question 2. Corollary 6.1 realizes
exponent \(1/3\); the all-exponents question remains posed. Searches for
the title, authors, and marginal-growth exponents through 2026 located
no complete answer. Infinite compact families and subsequence-only lower
bounds do not meet the finite-family, every-length target.

\subsection{Audit --- 2026-09-10}\label{audit-2026-09-10}

Rechecked \href{https://arxiv.org/html/2209.00449}{Varney--Morris,
Corollary 6.1, Proposition 3.1, and Question 2}: exponent \(1/3\) is
realized and realizable exponents are closed under addition. The
all-exponents assertion remains open. Author and
marginal-growth-exponent searches found no complete finite-family
construction or obstruction.
\end{document}
